\subsection{Metody ładowania danych}

Wspólnym punktem wyjścia dla wszystkich czterech silników są pliki CSV
wyprodukowane przez generator (\texttt{src/generators/data\_generator.py})
i zapisane w katalogu \texttt{app/data/\{small,medium,large\}/} -- po
jednym pliku na tabelę modelu relacyjnego. Każdy z czterech loaderów
(\texttt{postgres\_loader.py}, \texttt{mysql\_loader.py}, \texttt{mongo\_loader.py},
\texttt{neo4j\_loader.py}) czyta te same pliki niezależnie i przekształca je
do natywnej reprezentacji docelowego silnika. Dzięki temu cztery bazy
otrzymują identyczne dane wejściowe -- różni je wyłącznie sposób, w jaki
te dane zostają wsadowo wprowadzone do silnika.

W każdym z loaderów obowiązuje ta sama, świadomie przyjęta strategia
dwuetapowa: najpierw ładowane są wszystkie wiersze do pustych tabel
i kolekcji \emph{bez indeksów dodatkowych}, a dopiero w drugim kroku -- po
zakończeniu importu -- wywoływana jest osobna metoda \texttt{create\_indexes()}
budująca indeksy na komplecie danych. Klucze główne i ograniczenia
unikalności zdefiniowane w schemacie pozostają aktywne przez cały czas;
dotyczy to wyłącznie indeksów dodatkowych przyspieszających scenariusze SELECT.
Takie podejście jest standardową praktyką ładowania masowego: koszt
przebudowy struktur indeksowych po zakończeniu importu jest istotnie
niższy niż koszt utrzymywania ich w spójności po każdym pojedynczym
wstawieniu. Dodatkowo metoda ta jest wymagana przez strukturę pipeline'u
benchmarków -- pomiary scenariuszy CRUD są wykonywane dwukrotnie:
najpierw na bazie bez indeksów dodatkowych, a następnie po ich założeniu,
co pozwala bezpośrednio porównać wpływ indeksów na czas wykonania
zapytań.

\subsubsection*{PostgreSQL -- COPY FROM STDIN}

PostgreSQL ładowany jest mechanizmem \texttt{COPY FROM STDIN} -- najszybszą
natywną metodą importu masowego dostępną w tym silniku. W odróżnieniu
od ciągu \texttt{INSERT}-ów, \texttt{COPY} pomija znaczną część narzutu
parsera i protokołu komunikacyjnego, omija logikę triggerów per-statement
oraz pozwala silnikowi grupować zapis stron WAL. Loader korzysta
z driveira \texttt{psycopg~3} i jego strumieniowego API \texttt{cursor.copy()}
karmionego kolejnymi 8-megabajtowymi blokami pliku CSV. Przed właściwym
ładowaniem podnoszone są dwa parametry sesji:

\begin{lstlisting}[language=SQL,basicstyle=\ttfamily\small]
SET maintenance_work_mem = '512MB';
SET work_mem = '64MB';
\end{lstlisting}

\noindent \texttt{maintenance\_work\_mem} przyspiesza późniejsze tworzenie
indeksów (krok \texttt{CREATE INDEX} może wówczas posortować dane w pamięci
zamiast korzystać z plików tymczasowych), \texttt{work\_mem} zwiększa
budżet pamięciowy operacji sortujących i hash-ujących wykorzystywanych
przez \texttt{COPY} przy walidacji kluczy obcych. Tabele są ładowane
w deterministycznej kolejności respektującej zależności kluczy obcych
(\texttt{users}~$\rightarrow$~\texttt{profiles}~$\rightarrow$~\texttt{subscriptions}~$\rightarrow$~\texttt{payments}~$\rightarrow$~\dots),
a po zakończeniu importu sekwencje (\texttt{users\_user\_id\_seq} i analogiczne)
są resetowane funkcją \texttt{setval()} do wartości \texttt{MAX(pk)+1} -- co
zabezpiecza późniejsze scenariusze \texttt{INSERT} przed konfliktami kluczy
głównych.

\subsubsection*{MySQL -- LOAD DATA LOCAL INFILE}

MySQL wykorzystuje swój odpowiednik \texttt{COPY} -- instrukcję
\texttt{LOAD DATA LOCAL INFILE}, która wczytuje plik CSV bezpośrednio
po stronie serwera, omijając całkowicie warstwę przygotowywanych
zapytań. Aby mechanizm ten zadziałał z poziomu klienta, w pliku
\texttt{docker-compose.yml} ustawiona jest opcja serwera
\texttt{local-infile=1}, a w driverze PyMySQL flaga
\texttt{local\_infile=True} przekazywana do \texttt{connect()}.
Bezpośrednio przed importem loader wyłącza trzy mechanizmy
narzutowe:

\begin{lstlisting}[language=SQL,basicstyle=\ttfamily\small]
SET FOREIGN_KEY_CHECKS = 0;
SET UNIQUE_CHECKS = 0;
SET AUTOCOMMIT = 0;
\end{lstlisting}

\noindent Wyłączenie \texttt{FOREIGN\_KEY\_CHECKS} pozwala załadować
tabele potomne, zanim wszystkie wiersze tabeli rodzica zostaną
przeniesione na dysk; \texttt{UNIQUE\_CHECKS} zatrzymuje tymczasową
walidację unikalności (klucze i tak są spójne, ponieważ generator
gwarantuje brak duplikatów); wyłączenie \texttt{AUTOCOMMIT} grupuje
zapisy w jedną transakcję, redukując liczbę synchronizacji pliku
\texttt{redo log}. Po zakończeniu importu wszystkie trzy flagi są
przywracane, a transakcja zatwierdzona pojedynczym \texttt{COMMIT}.
Dla tabel, w których oryginalny CSV reprezentuje wartości logiczne
jako napisy \texttt{true}/\texttt{false} (\texttt{profiles},
\texttt{subscriptions}, \texttt{content}, \texttt{watch\_history}),
loader generuje wzbogaconą postać \texttt{LOAD DATA} ze zmiennymi
sesyjnymi i klauzulą \texttt{SET}, wykonującymi konwersję
napis~$\rightarrow$~\texttt{TINYINT} po stronie serwera.

\subsubsection*{MongoDB -- insert\_many w trybie nieuporządkowanym}

MongoDB nie posiada równoważnika \texttt{COPY} -- ładowanie odbywa się
metodą \texttt{insert\_many()} drivera PyMongo, której PyMongo automatycznie
serializuje wsady do protokołu \texttt{OP\_MSG}. Kluczowe są dwie decyzje
projektowe. Po pierwsze, dokumenty są wkładane wsadami po
\textbf{25\,000 pozycji} (\texttt{BATCH\_SIZE = 25000}); jest to wartość
będąca kompromisem -- mniejsze paczki podnoszą udział narzutu sieciowego,
większe ryzykują przekroczenie domyślnego limitu wiadomości BSON
(48~MB) dla tabeli \texttt{content}, w której każdy dokument zawiera
zagnieżdżoną listę sezonów, odcinków i obsady. Po drugie, każde
wywołanie \texttt{insert\_many()} korzysta z flagi \texttt{ordered=False}.
W trybie nieuporządkowanym MongoDB wykonuje wstawiania równolegle
i nie przerywa wsadu po pierwszym błędzie -- co istotnie zwiększa
przepustowość i odporność na pojedyncze duplikaty.

Specyfika modelu dokumentowego wymaga, by przed właściwym \texttt{insert\_many}
loader scalił dane z wielu plików CSV w jeden dokument. Kolekcja
\texttt{users} powstaje z połączenia \texttt{users.csv}, \texttt{profiles.csv}
oraz aktywnych rekordów z \texttt{subscriptions.csv} osadzonych
bezpośrednio w polu \texttt{subscription}. Kolekcja \texttt{content}
absorbuje dane z \texttt{content.csv}, \texttt{content\_people.csv},
\texttt{seasons.csv} oraz \texttt{episodes.csv} budując w pamięci
hierarchię \texttt{content}~$\rightarrow$~\texttt{seasons[]}~$\rightarrow$~\texttt{episodes[]}.
Wsadowy import po stronie sterownika zaczyna się dopiero po skompletowaniu
całej struktury w słowniku Pythona, co dla wolumenu \emph{large} oznacza
przejściowy szczyt zużycia RAM warstwy orkiestracji rzędu kilkuset MB --
celowo akceptowalny, ponieważ alternatywa (wielokrotne aktualizacje
\texttt{\$push} po stronie bazy) byłaby o rzędy wielkości wolniejsza.

\subsubsection*{Neo4j -- UNWIND w sesjach Bolt}

Neo4j jest najtrudniejszym przypadkiem, ponieważ ładowanie dotyczy
nie tylko węzłów, ale również relacji. Loader korzysta z konstrukcji
Cypher \texttt{UNWIND \$rows AS r CREATE \dots} przyjmującej
parametr \texttt{rows} jako tablicę słowników -- zamiast wysyłać
N pojedynczych zapytań \texttt{CREATE}, wykonujemy jedno zapytanie
przetwarzające cały wsad. Również tutaj wsad ma rozmiar
\textbf{25\,000} elementów; powyżej tej wielkości pojawia się ryzyko
przekroczenia budżetu transakcji (\texttt{dbms.memory.transaction.total.max
~=~3g}) dla rzędu wielkości scenariusza \emph{large}.

Ładowanie dzieli się na dwa wyraźne etapy. W pierwszym tworzone
są \emph{ograniczenia unikalności} (\texttt{CREATE CONSTRAINT \dots
REQUIRE x IS UNIQUE}) na identyfikatorach kluczowych etykiet --
\texttt{User}, \texttt{Profile}, \texttt{Content}, \texttt{Season},
\texttt{Episode}, \texttt{Person}, \texttt{Genre},
\texttt{Subscription} i \texttt{Payment}. W Neo4j ograniczenie
unikalności jest jednocześnie \emph{indeksem zakresowym} -- stworzenie
go przed importem jest absolutnie kluczowe, ponieważ druga faza
(tworzenie relacji) wykonuje miliony operacji
\texttt{MATCH~(n~\{id:~r.id\})}, które bez indeksu degenerowałyby
się do skanu wszystkich węzłów danej etykiety. W drugim etapie
ładowane są węzły wszystkich etykiet, a następnie -- w równie
deterministycznej kolejności -- relacje (\texttt{HAS\_PROFILE},
\texttt{HAS\_SUBSCRIPTION}, \texttt{HAS\_PAYMENT}, \texttt{HAS\_GENRE},
\texttt{HAS\_SEASON}, \texttt{HAS\_EPISODE}, \texttt{ACTED\_IN},
\texttt{DIRECTED}, \texttt{WROTE}, \texttt{WATCHED},
\texttt{ADDED\_TO\_LIST}, \texttt{RATED}). Plik \texttt{watch\_history.csv}
dla wolumenu \emph{large} liczy miliony wierszy -- aby uniknąć
materializacji całej listy w pamięci, dla relacji \texttt{WATCHED},
\texttt{ADDED\_TO\_LIST} oraz \texttt{RATED} loader korzysta
z generatorów Pythona i metody \texttt{\_batch\_write\_streaming()},
która buforuje paczkę w pamięci tylko do osiągnięcia \texttt{BATCH\_SIZE},
po czym natychmiast oddaje ją do bazy.

\subsubsection*{Czyszczenie środowiska przed kolejnym ładowaniem}

W pipeline benchmarków każda baza może być ładowana wielokrotnie --
różnymi wolumenami danych oraz w wariancie z indeksami i bez. Aby
uniknąć pozostałości po poprzednich uruchomieniach, każdy loader
zaczyna od czyszczenia. PostgreSQL i MySQL korzystają z \texttt{TRUNCATE}
(odpowiednio z opcją \texttt{CASCADE} oraz po wyłączeniu kontroli
kluczy obcych), MongoDB wykonuje \texttt{drop\_collection()}
na wszystkich kolekcjach docelowych, a Neo4j wykonuje pętlę
\texttt{MATCH~()-[r]->() WITH r LIMIT 50000 DELETE r} oraz
\texttt{MATCH~(n) WITH n LIMIT 50000 DELETE n}, dopóki licznik
usuniętych elementów nie spadnie do zera. Wsadowe usuwanie zamiast
pojedynczego \texttt{MATCH (n) DETACH DELETE n} jest podyktowane
limitem pamięci transakcyjnej silnika -- jedna transakcja usuwająca
kilkanaście milionów relacji wyczerpałaby budżet 3~GB.

% -------------------------------------------------------
\subsection{Typy indeksów i strategia indeksowania}

Drugim etapem każdego ładowania -- wykonywanym przez metodę
\texttt{create\_indexes()} dopiero po zakończeniu importu -- jest
budowa zestawu indeksów dodatkowych. Dobór indeksów został podporządkowany
24~scenariuszom benchmarków: każdy istotny predykat \texttt{WHERE},
warunek \texttt{JOIN} oraz klauzula \texttt{ORDER BY} z gorącej ścieżki
SELECT-ów ma wsparcie w postaci dedykowanego indeksu, co pozwala
hipotezie pracy wybrzmieć w pomiarach. Jednocześnie zakres
indeksów jest świadomie ograniczony do struktur o praktycznym
zastosowaniu -- nie definiujemy indeksów martwych, których jedyną
funkcją byłoby spowolnienie operacji zapisu.

Wszystkie cztery silniki dzielą zestaw indeksów ,,bazowych'' --
B-drzewa pojedynczych kolumn będących kluczami obcymi lub
najczęstszymi predykatami filtrowania (\texttt{users.status},
\texttt{content.type}, \texttt{watch\_history.content\_id} i~analogiczne).
Powyżej tego wspólnego mianownika każdy z silników eksponuje
typy indeksów charakterystyczne dla swojej architektury -- i to
właśnie te ,,specjalistyczne'' struktury przesądzają o sile
poszczególnych baz w wybranych scenariuszach.

\subsubsection*{PostgreSQL}

Loader \texttt{postgres\_loader.py} tworzy 17 indeksów dodatkowych,
spośród których pięć wykracza poza standardowe B-drzewo:

\begin{itemize}
    \item \textbf{Indeksy złożone (composite B-tree)} --
    \texttt{idx\_subscriptions\_status\_plan} po
    \texttt{(status, plan\_name)} oraz
    \texttt{idx\_wh\_profile\_started} po
    \texttt{(profile\_id, started\_at DESC)}. Drugi obsługuje
    najczęstsze zapytanie domeny aktywności -- ,,ostatnie sesje
    profilu'' -- pozwalając jednym przejściem indeksu odczytać
    posortowany wynik bez kroku \texttt{Sort}.
    \item \textbf{Indeks częściowy (partial index)} --
    \texttt{idx\_content\_active\_popular} założony
    na \texttt{(popularity\_score DESC)} z klauzulą
    \texttt{WHERE is\_active = TRUE}. Indeks fizycznie
    przechowuje tylko aktywne tytuły, co dla katalogu z
    nieaktywnymi pozycjami daje istotnie mniejszą strukturę
    i szybszy dostęp w scenariuszach typu ,,top~N popularnych
    treści''.
    \item \textbf{GIN z trygramami} --
    \texttt{idx\_content\_title\_trgm} na kolumnie
    \texttt{title} z operatorem \texttt{gin\_trgm\_ops}
    z rozszerzenia \texttt{pg\_trgm} (loader włącza je
    instrukcją \texttt{CREATE EXTENSION IF NOT EXISTS pg\_trgm}).
    Indeks ten przyspiesza zapytania z operatorem
    \texttt{LIKE '\%fragment\%'} -- bez niego predykat z wiodącym
    znakiem wieloznacznym wymusiłby \texttt{Seq Scan}.
    \item \textbf{GIN dla JSONB} --
    \texttt{idx\_content\_metadata\_gin} na kolumnie
    \texttt{metadata} typu \texttt{jsonb}. Daje dostęp
    indeksowy do operatorów ,,zawiera'' \texttt{@>}
    i ,,istnieje klucz'' \texttt{?} -- kluczowe dla scenariuszy
    odpytujących pola zagnieżdżone (np. nazwę studia
    produkcyjnego z \texttt{metadata->>'studio'}).
\end{itemize}

\subsubsection*{MySQL}

Loader \texttt{mysql\_loader.py} tworzy 10 indeksów dodatkowych
o nieco węższym profilu, wynikającym z ograniczeń
silnika InnoDB:

\begin{itemize}
    \item \textbf{Indeks złożony} --
    \texttt{idx\_wh\_profile\_started} po
    \texttt{(profile\_id, started\_at DESC)} -- analogicznie
    do PostgreSQL.
    \item \textbf{FULLTEXT} -- \texttt{idx\_content\_title\_ft}
    na kolumnie \texttt{title}. MySQL nie udostępnia
    odpowiednika trygramowego GIN, dlatego dla scenariuszy
    z wyszukiwaniem po fragmencie tytułu wykorzystywany jest
    natywny mechanizm pełnotekstowy z operatorem
    \texttt{MATCH \dots AGAINST}.
    \item \textbf{Indeks funkcyjny po polu JSON} --
    \texttt{idx\_content\_metadata\_studio} założony na
    wyrażeniu \texttt{(CAST(metadata->>'\$.studio' AS CHAR(100)))}.
    InnoDB nie indeksuje całych dokumentów JSON tak jak
    PostgreSQL przez GIN -- zamiast tego pozwala zaindeksować
    pojedyncze, jawnie wyciągnięte pole. Wybór został
    podyktowany scenariuszami benchmarka, które filtrują po
    nazwie studia.
\end{itemize}

\subsubsection*{MongoDB}

W modelu dokumentowym indeksy są zakładane na poziomie
kolekcji metodą \texttt{create\_index()} sterownika PyMongo.
Loader buduje ich 13:

\begin{itemize}
    \item \textbf{Indeksy pojedynczych pól} po
    \texttt{users.email} (\texttt{unique=True}),
    \texttt{users.status}, \texttt{content.type},
    \texttt{content.is\_active} i analogicznych.
    \item \textbf{Indeks wielokluczowy (multikey)} po polu
    \texttt{content.genres} -- pole jest tablicą napisów,
    więc MongoDB automatycznie tworzy indeks wielokluczowy
    pozwalający efektywnie odpytywać po wartości jednego
    z elementów (\texttt{\{ genres: "drama" \}}).
    \item \textbf{Indeks zagnieżdżony} po
    \texttt{users.profiles.profile\_id} -- ,,profil''
    jest osadzony w dokumencie użytkownika, a indeks
    pozwala szybko odnaleźć dokument-rodzic po identyfikatorze
    osadzonego elementu.
    \item \textbf{Indeks złożony} po
    \texttt{watch\_history (profile\_id ASC, started\_at DESC)}
    -- kolejność identyczna jak w PostgreSQL i MySQL.
    \item \textbf{Indeks tekstowy} typu \texttt{text} na
    polu \texttt{content.title} -- obsługuje natywny operator
    \texttt{\$text} dla zapytań wyszukiwania pełnotekstowego.
    \item \textbf{Indeksy z ograniczeniem unikalności} po
    \texttt{(profile\_id, content\_id)} dla kolekcji
    \texttt{ratings} oraz \texttt{my\_list} -- jednocześnie
    zabezpieczają model przed duplikatami i przyspieszają
    najczęstszą formę dostępu do tych kolekcji.
\end{itemize}

\subsubsection*{Neo4j}

Neo4j od najwyższego poziomu rozróżnia kilka typów indeksów,
a dodatkowo każde ograniczenie unikalności automatycznie tworzy
indeks zakresowy. Loader \texttt{neo4j\_loader.py} korzysta
z trzech kategorii:

\begin{itemize}
    \item \textbf{Indeksy zakresowe (range index)} -- domyślny
    typ tworzony przez \texttt{CREATE INDEX FOR (n:Label)
    ON (n.property)}. Pokrywa pola \texttt{User.status},
    \texttt{User.country\_code}, \texttt{Content.type},
    \texttt{Content.popularity\_score}, \texttt{Content.is\_active},
    \texttt{Subscription.status} i \texttt{Payment.status}.
    Dodatkowo, \emph{niejawnie}, indeksy zakresowe powstają
    dla wszystkich identyfikatorów objętych ograniczeniami
    unikalności tworzonymi w pierwszej fazie ładowania
    (\texttt{User.user\_id}, \texttt{Content.content\_id}
    i osiem analogicznych).
    \item \textbf{Indeksy tekstowe (text index)} -- pokrywają
    pola \texttt{Content.title} oraz \texttt{Content.metadata}.
    Indeks tekstowy w Neo4j jest zoptymalizowany pod operatory
    \texttt{STARTS WITH}, \texttt{ENDS WITH} i \texttt{CONTAINS}
    -- czyli te same przypadki, w których PostgreSQL korzysta
    z trygramowego GIN.
    \item \textbf{Indeks pełnotekstowy (fulltext index)}
    \texttt{content\_title\_ft} -- używany w scenariuszach
    odpytywanych przez procedurę \texttt{db.index.fulltext.queryNodes()}.
    Indeks pełnotekstowy oparty jest na bibliotece Lucene
    i obsługuje analizator językowy oraz scoring -- czego nie
    daje samodzielny indeks tekstowy.
\end{itemize}

\subsubsection*{Strategia tworzenia i kasowania indeksów w pipeline'ie}

Cały pipeline benchmarków został zaprojektowany tak, aby tę samą
bazę odpytać w dwóch wariantach -- bez i z indeksami dodatkowymi --
przy minimalnym narzucie organizacyjnym. Każdy loader udostępnia
parę symetrycznych metod: \texttt{create\_indexes()} oraz
\texttt{drop\_indexes()}. Pierwsza wykonuje listę instrukcji
\texttt{CREATE INDEX IF NOT EXISTS}, druga -- listę
\texttt{DROP INDEX IF EXISTS} dla każdego indeksu z osobna
po nazwie. W przypadku Neo4j, w którym lista indeksów
nieskojarzonych z ograniczeniami zmienia się w czasie życia
bazy, \texttt{drop\_indexes()} odpytuje katalog systemowy
zapytaniem \texttt{SHOW INDEXES YIELD name, owningConstraint
WHERE owningConstraint IS NULL} i kasuje zwrócone pozycje --
co gwarantuje, że indeksy podtrzymujące unikalność identyfikatorów
nie zostaną usunięte.

Komenda \texttt{run-all} pliku \texttt{main.py} sekwencyjnie
wywołuje cały cykl: \emph{generate}~$\rightarrow$~\emph{load}
(bez indeksów dodatkowych)~$\rightarrow$~\emph{benchmark}~$\rightarrow$
\emph{explain}~$\rightarrow$~\emph{create-indexes}~$\rightarrow$~\emph{benchmark}
(z indeksami)~$\rightarrow$~\emph{explain}~$\rightarrow$~\emph{visualize}.
Tym samym wynik każdego scenariusza dostępny jest w dwóch wariantach
opisanych jednoznaczną kolumną \texttt{indexed} w pliku CSV
z wynikami, a późniejsza analiza (mapy cieplne, weryfikacja
hipotezy w sekcji 11) sprowadza się do wyliczenia stosunku
czasów \texttt{indexed} / \texttt{not\_indexed} dla każdego
scenariusza i każdej bazy.
